% =====================================================================
%  Estudio de casos --- MDS221 Inteligencia Artificial
%  Basado en clase_1_4.tex (¿Por qué funciona el aprendizaje profundo?)
%  Dataset: Maternal Health Risk (UCI id 863) vía Tzimourta et al. (2025)
% =====================================================================
\documentclass[11pt, a4paper]{article}

\usepackage[utf8]{inputenc}
\usepackage[T1]{fontenc}
\usepackage[provide=*,spanish]{babel}
\usepackage[margin=2.1cm]{geometry}
\usepackage{amsmath, amssymb, bm}
\usepackage{booktabs}
\usepackage{tabularx}
\usepackage{xcolor}
\usepackage{enumitem}
\usepackage{tcolorbox}
\tcbuselibrary{skins, breakable}
\usepackage{titlesec}
\usepackage[colorlinks=true, allbordercolors=white]{hyperref}

% ---------- Paleta UCM (la misma de las clases) ----------
\definecolor{UCMBlue}{RGB}{4, 171, 226}
\definecolor{UCMNavy}{RGB}{20, 65, 134}
\definecolor{AccentGreen}{RGB}{0, 130, 100}
\definecolor{StanfordRed}{RGB}{140, 21, 21}
\definecolor{StanfordGray}{RGB}{77, 77, 77}
\definecolor{LightGray}{RGB}{240, 240, 240}

\hypersetup{linkcolor=UCMNavy, urlcolor=UCMNavy, citecolor=UCMNavy}

\titleformat{\section}{\normalfont\Large\bfseries\color{UCMNavy}}{\thesection.}{0.6em}{}
\titleformat{\subsection}{\normalfont\large\bfseries\color{UCMBlue!80!black}}{\thesubsection}{0.6em}{}
\titlespacing*{\section}{0pt}{1.4em}{0.6em}

\setlist[itemize]{leftmargin=1.2em, itemsep=2pt, topsep=3pt}
\setlist[enumerate]{leftmargin=1.4em, itemsep=2pt, topsep=3pt}

% ---------- Cajas ----------
\newtcolorbox{pregunta}{colback=UCMBlue!7, colframe=UCMBlue, arc=3pt,
	left=8pt, right=8pt, top=5pt, bottom=5pt, boxrule=0.9pt}
\newtcolorbox{aviso}{colback=StanfordRed!5, colframe=StanfordRed, arc=3pt,
	left=8pt, right=8pt, top=5pt, bottom=5pt, boxrule=0.9pt}
\newtcolorbox{entrega}{colback=AccentGreen!6, colframe=AccentGreen, arc=3pt,
	left=8pt, right=8pt, top=5pt, bottom=5pt, boxrule=0.9pt}

% Vínculo explícito con la clase
\newcommand{\enlace}[1]{\par\vspace{2pt}{\footnotesize\color{StanfordGray}%
	$\hookrightarrow$\  #1}}

\begin{document}

% =====================================================================
\begin{center}
{\LARGE\bfseries\color{UCMNavy} Estudio de casos}\\[4pt]
{\Large\bfseries Aprendizaje profundo en la detección de riesgo materno}\\[10pt]
{\normalsize MDS221 --- Inteligencia Artificial \quad$\bullet$\quad Unidad 1}\\[2pt]
{\normalsize Mag\'ister en Data Science --- Universidad Cat\'olica del Maule}\\[2pt]
{\small Sergio Hern\'andez, PhD \quad$\bullet$\quad \texttt{shernandez@ucm.cl}}
\end{center}
\vspace{0.4em}
\hrule
\vspace{1em}

\noindent
\begin{tabularx}{\textwidth}{@{}p{3.2cm}X@{}}
\textbf{Modalidad} & Individual.\\
\textbf{Dedicaci\'on} & Aproximadamente 12 horas de trabajo aut\'onomo.\\
\textbf{Entregables} & Notebook reproducible \texttt{.ipynb} + informe en PDF (6--8 p\'aginas).\\
\textbf{Evaluaci\'on} & R\'ubrica de la secci\'on~\ref{sec:rubrica} (240 puntos, 60\,\% de exigencia).\\
\end{tabularx}

\vspace{0.8em}
\begin{pregunta}
\textbf{El caso en una frase.} Un art\'iculo publicado en \emph{Healthcare} (2025) compara seis clasificadores para predecir el nivel de riesgo de una embarazada a partir de seis signos vitales. El \textbf{bosque aleatorio alcanza 84{,}5\,\%} de exactitud; las \textbf{redes neuronales se quedan en 69\,\%}. La sobreparametrizaci\'on y la profundidad son justamente lo que hace funcionar al aprendizaje profundo. \textbf{¿Por qu\'e aqu\'i pierde?}
\end{pregunta}

% =====================================================================
\section{Contexto clínico}

La mortalidad materna sigue siendo uno de los indicadores de salud m\'as sensibles a la desigualdad: la mayor\'ia de las muertes ocurre en contextos con acceso tard\'io o intermitente al control prenatal, y una fracci\'on importante es evitable si el riesgo se detecta a tiempo. En ese escenario, un sistema de \emph{triage} que priorice a qui\'en derivar --- construido sobre mediciones baratas y repetibles (presi\'on, glicemia, temperatura, pulso) --- tiene un valor cl\'inico directo, incluso si su exactitud es modesta.

Tzimourta et al.~(2025) exploran precisamente eso: entrenan seis clasificadores sobre registros recolectados mediante un sistema de monitoreo IoT en hospitales, clínicas comunitarias y centros de salud materna de Bangladesh, y comparan su desempe\~no para clasificar el riesgo en tres niveles.

% =====================================================================
\section{El conjunto de datos}

\noindent\textbf{Maternal Health Risk} --- UCI Machine Learning Repository, id 863.
DOI \texttt{10.24432/C5DP5D}, licencia CC BY 4.0.\\
\url{https://archive.ics.uci.edu/dataset/863/maternal+health+risk}

\vspace{0.6em}
\noindent
\begin{tabularx}{\textwidth}{@{}llX@{}}
\toprule
\textbf{Variable} & \textbf{Unidad} & \textbf{Descripci\'on}\\
\midrule
\texttt{Age}          & a\~nos      & Edad de la gestante.\\
\texttt{SystolicBP}   & mmHg        & Presi\'on arterial sist\'olica.\\
\texttt{DiastolicBP}  & mmHg        & Presi\'on arterial diast\'olica.\\
\texttt{BS}           & mmol/L      & Glicemia (\emph{blood sugar}).\\
\texttt{BodyTemp}     & $^\circ$F   & Temperatura corporal.\\
\texttt{HeartRate}    & lpm         & Frecuencia card\'iaca en reposo.\\
\midrule
\texttt{RiskLevel}    & ---         & \textbf{Objetivo}: \emph{low risk} (406), \emph{mid risk} (336), \emph{high risk} (272).\\
\bottomrule
\end{tabularx}

\vspace{0.6em}
\noindent
Seis atributos, tres clases, y del orden de mil registros. Se descarga con:

\vspace{0.3em}
\noindent\colorbox{LightGray}{\parbox{\dimexpr\textwidth-2\fboxsep}{\ttfamily\small
from ucimlrepo import fetch\_ucirepo\\
mhr = fetch\_ucirepo(id=863)\\
X, y = mhr.data.features, mhr.data.targets
}}

\begin{comment}
\textbf{Primera discrepancia, y es suya para resolver.} El art\'iculo declara \textbf{1014} instancias (1013 tras corregir dos frecuencias card\'iacas imposibles de 7\,lpm); la ficha actual de UCI reporta \textbf{1013}. Verifique cu\'antas filas trae su descarga, cu\'antas est\'an \textbf{duplicadas exactamente} y d\'ejelo documentado antes de entrenar nada. La cifra que obtenga condiciona todo lo que sigue.
\end{comment}

% =====================================================================
\section{Linea Base}

Los autores entrenan con validaci\'on cruzada de 10 pliegues y reportan:

\vspace{0.5em}
\noindent
\begin{tabularx}{\textwidth}{@{}Xccc@{}}
\toprule
\textbf{Clasificador} & \textbf{Exactitud} & \textbf{Precisi\'on} & \textbf{Sensibilidad (TP rate)}\\
\midrule
Na\"ive Bayes                     & 59{,}07\,\% & 58{,}10\,\% & 59{,}10\,\%\\
Perceptr\'on multicapa (MLP)      & 69{,}03\,\% & 61{,}04\,\% & 61{,}01\,\%\\
Red totalmente conectada (FCNN)   & 68{,}93\,\% & 68{,}96\,\% & 68{,}91\,\%\\
M\'aquina de soporte vectorial    & 75{,}74\,\% & 76{,}30\,\% & 75{,}70\,\%\\
\'Arbol de decisi\'on             & 78{,}60\,\% & 79{,}00\,\% & 78{,}60\,\%\\
\textbf{Bosque aleatorio}         & \textbf{84{,}52\,\%} & \textbf{84{,}93\,\%} & \textbf{84{,}52\,\%}\\
\textbf{Bosque aleatorio + SMOTE} & \textbf{88{,}03\,\%} & \textbf{88{,}10\,\%} & \textbf{88{,}00\,\%}\\
\bottomrule
\end{tabularx}

\vspace{0.8em}
\noindent
Las dos arquitecturas neuronales quedan \textbf{\'ultimas entre los m\'etodos no triviales}, quince puntos por debajo del bosque. Ese es el hecho a explicar. Los mecanismos que explican el desempeño del Aprendizaje Profundo son : sobreparametrizaci\'on, umbral de interpolaci\'on, doble descenso, planitud del m\'inimo, sesgo inductivo, billete de loter\'ia. El trabajo consiste en \textbf{usar ese vocabulario como instrumento de medici\'on}, no como decorado.

\begin{pregunta}
\textbf{Pregunta central del caso.} ¿El fracaso relativo de las redes neuronales en este problema es (a) un problema de \emph{capacidad} --- faltan par\'ametros o datos ---, (b) un problema de \emph{optimizaci\'on} --- la soluci\'on existe pero el entrenamiento no la encuentra ---, o (c) un problema de \emph{sesgo inductivo} --- la arquitectura no encaja con la estructura del dato? Su informe debe defender una respuesta \textbf{con evidencia propia}, no con citas.
\end{pregunta}

% =====================================================================
\section{Tareas}

Cada tarea indica con qu\'e parte de la clase 1.4 se conecta. Todos los experimentos se reportan con \textbf{al menos 5 semillas} y con media $\pm$ desviaci\'on est\'andar: una diferencia de dos puntos sin barras de error no es un resultado.

\subsection*{T1 --- Línea base y auditoría del dato}
\begin{enumerate}
	\item Reproduzca la l\'inea base: bosque aleatorio y MLP con validaci\'on cruzada estratificada de 10 pliegues. Compare sus n\'umeros con la tabla anterior y explique cualquier brecha mayor a 3 puntos.
	\item Cuente las filas duplicadas exactas. Reeval\'ue asegurando que \textbf{ning\'un duplicado quede a ambos lados} de la partici\'on. ¿Cu\'anta exactitud se pierde? Discuta qu\'e parte del 84{,}5\,\% original era memorizaci\'on medida sobre s\'i misma.
	\item Con clases 406/336/272, argumente por qu\'e la exactitud sola es un mal resumen. Reporte matriz de confusi\'on y F1 por clase.
\end{enumerate}
%\enlace{secci\'on ``¿Son los datos?'' --- la advertencia de Zhang et al.\ es que una red memoriza \emph{cualquier} conjunto finito; su validaci\'on tiene que estar dise\~nada para notar la diferencia.}

\subsection*{T2 --- Generalización y sobreparametrización}
\begin{enumerate}
	\item Permute \texttt{RiskLevel} al azar, destruyendo toda relaci\'on con las variables. Entrene MLP de ancho creciente hasta llevar el error de \textbf{entrenamiento} a cero.
	\item Reporte el n\'umero m\'inimo de par\'ametros necesario y comp\'arelo con $n \approx 1000$ ejemplos. ¿Cu\'antas veces sobreparametrizado quedó el modelo?
	\item Repita con las etiquetas reales. ¿Necesita m\'as o menos capacidad? ¿Qu\'e dice eso sobre la ``simplicidad'' de la funci\'on subyacente?
\end{enumerate}
%\enlace{r\'eplica en miniatura del experimento de Zhang et al.\ (2017) que aparece en la l\'amina de p\'ixeles y etiquetas %aleatorias.}

\subsection*{T3 --- Dilema Sesgo-Varianza}
\begin{enumerate}
	\item MLP de una capa oculta con ancho en escala logar\'itmica, de 2 a 2048 unidades. Grafique error de entrenamiento y de prueba contra n\'umero de par\'ametros.
	\item Marque el \textbf{umbral de interpolaci\'on} (par\'ametros $\approx$ datos). ¿Observa el pico y el segundo descenso, o solo la curva en U cl\'asica?
	\item Qu\'e condiciones no se cumplen aqu\'i (tama\~no de $n$, ruido de etiqueta, regularizaci\'on impl\'icita del optimizador).
\end{enumerate}
%\enlace{l\'amina ``Sobreparametrizaci\'on y doble descenso''. La curva en U cl\'asica no est\'a mal, est\'a \emph{incompleta} %--- pero solo cuando se llega al r\'egimen moderno.}

\subsection*{T4 --- Hiperpar\'ametros}
\begin{enumerate}
	\item Grilla de tama\~no de lote $\in \{8, 32, 128, 512\}$ por tasa de aprendizaje $\in \{10^{-3}, 10^{-2}, 10^{-1}\}$. Grafique el error de prueba contra la \textbf{raz\'on lote/tasa}.
	\item Estime la planitud de cada m\'inimo: p\'erdida m\'axima en una vecindad $\lVert\bm\delta\rVert = \varepsilon$ alrededor de los par\'ametros finales (criterio de Keskar), con al menos 20 direcciones aleatorias.
	\item ¿Correlacionan planitud y generalizaci\'on \textbf{en este problema}? Contraste su resultado con la advertencia de Dinh et al.\ sobre reparametrizaciones triviales de ReLU.
\end{enumerate}
%\enlace{l\'aminas ``El algoritmo elige a cu\'al m\'inimo llegar'' y ``Planitud del m\'inimo''.}

\subsection*{T5 --- Sesgo inductivo}
\begin{enumerate}
	\item Formule una hip\'otesis concreta sobre por qu\'e un \emph{ensemble} de \'arboles supera a una red en datos tabulares peque\~nos, y dise\~ne \textbf{un experimento que pueda refutarla}.
	\item Contraprueba obligatoria: intente cerrar la brecha con al menos dos intervenciones (normalizaci\'on por variable, regularizaci\'on fuerte, \emph{early stopping}, \emph{embeddings} de variables discretizadas, \emph{ensemble} de redes). Reporte cu\'anto se acorta.
	\item Relacione su conclusi\'on con las tres explicaciones que la clase da para la profundidad: complejidad de la funci\'on, tratabilidad del entrenamiento y sesgo inductivo.
\end{enumerate}
%\enlace{l\'aminas ``Arquitectura: el sesgo inductivo'' y ``Tres explicaciones posibles para la profundidad''.}

\subsection*{T6 --- Modelos Fundacionales}
\begin{enumerate}
	\item Utilice el modelo fundacional TabFM\footnote{\url{https://github.com/google-research/tabfm/blob/main/examples/tabarena_classification_example.py}} para resolver el problema de clasficación usando promediado de logits (logit averaging).
	\item Utilice el modelo fundacional TabFM para resolver el problema de clasficación usando feature engineering (feature crosses / SVD features + NNLS-weighted blending + probability calibration).
	\item Compare los resultados de ambos enfoques con respecto a las m\'etricas de evaluaci\'on ROC, AUC y log-loss con respecto al mejor MLP obtenido.
\end{enumerate}

\vspace{0.3em}
\noindent\colorbox{LightGray}{\parbox{\dimexpr\textwidth-2\fboxsep}{\ttfamily\small
model = tabfm.tabfm\_v1\_0\_0\_pytorch.load(model\_type="classification")\\
results = \{\}\\
  results["default"] = \_evaluate(\\
      tabfm.TabFMClassifier(model=model, random\_state=SEED),\\
      x\_train, y\_train, x\_test, y\_test,\\
  )\\
  results["ensemble"] = \_evaluate(\\
      tabfm.TabFMClassifier.ensemble(model=model, random\_state=SEED),\\
      x\_train, y\_train, x\_test, y\_test,\\
  )\\
}}
%\enlace{l\'aminas ``Arquitectura: el sesgo inductivo'' y ``Tres explicaciones posibles para la profundidad''.}

\subsection*{T7 --- Lectura clínica: el error que no cuesta lo mismo}
\begin{enumerate}
	\item Un falso negativo de \emph{high risk} y un falso positivo de \emph{low risk} no son equivalentes. Proponga una matriz de costos justificada cl\'inicamente y reeval\'ue \textbf{todos} los modelos bajo ese criterio. ¿Cambia el ranking?
	\item Reporte expl\'icitamente la sensibilidad de la clase \emph{high risk} y ajuste umbrales o pesos de clase para maximizarla a costo controlado.
	\item Robustez: perturbe las se\~nales $\pm 5\,\%$ (error plausible de un tensi\'ometro o un term\'ometro) y mida la ca\'ida de desempe\~no de cada modelo. Discuta la analog\'ia con los ejemplos adversarios de la clase: ¿qu\'e tan lejos de la variedad de datos hay que ir para romper la predicci\'on?
	\item Los datos provienen de centros de salud de Bangladesh mediante sensores IoT. Discuta qu\'e supuestos habr\'ia que revisar antes de siquiera considerar un despliegue en la red asistencial chilena.
\end{enumerate}
%\enlace{l\'aminas ``Salir de la variedad de datos'' y ``Lo que todav\'ia no sabemos''.}

% =====================================================================
\section{Entregables}

\begin{entrega}
\begin{enumerate}[leftmargin=1.6em]
	\item \textbf{Notebook reproducible} (\texttt{.ipynb}). Debe correr de principio a fin sin intervenci\'on manual, con semillas fijadas y la descarga del dato incluida. Un cuaderno que no corre se eval\'ua como si no se hubiera entregado.
	\item \textbf{Informe} en PDF, 6 a 8 p\'aginas, con esta estructura m\'inima:
	\begin{itemize}
		\item Auditor\'ia del dato y l\'inea base corregida (T1).
		\item Resultados de T2--T5, cada uno con su figura y su lectura.
		\item Su respuesta a la pregunta central: (a) capacidad, (b) optimizaci\'on o (c) sesgo inductivo --- \textbf{y por qu\'e descarta las otras dos}.
		\item Secci\'on cl\'inica y \'etica (T7).
		\item Una secci\'on final de \textbf{l\'imites}: qu\'e no logr\'o establecer y qu\'e experimento habr\'ia hecho falta.
	\end{itemize}
	\item \textbf{Tabla resumen} de una p\'agina con todos los modelos entrenados: arquitectura, n\'umero de par\'ametros, protocolo, m\'etrica, media $\pm$ desviaci\'on.
\end{enumerate}
\end{entrega}

\vspace{0.4em}
\noindent
\textbf{Sobre el uso de asistentes de IA.} Est\'a permitido y se espera que lo declaren: incluyan un p\'arrafo indicando qu\'e herramienta usaron y para qu\'e. Lo que se eval\'ua es el criterio con que interpretan los resultados, y eso no se delega. Un informe con conclusiones que el propio cuaderno no respalda se penaliza en \emph{Precisi\'on y rigor}, independientemente de qui\'en lo haya escrito.

\begin{aviso}
\textbf{Un resultado negativo bien establecido vale tanto como uno positivo.} Si tras un barrido honesto la red sigue perdiendo contra el bosque, esa \emph{es} la conclusi\'on: la sobreparametrizaci\'on ayuda \textbf{a la escala y complejidad de los conjuntos que se usan hoy en visi\'on y lenguaje}, no que una red gane siempre.
\end{aviso}

% =====================================================================
\section{Rúbrica de evaluación}
\label{sec:rubrica}

\noindent
Seis criterios, cuatro niveles. Puntaje m\'aximo \textbf{240 puntos}.

\vspace{0.6em}
\footnotesize
\noindent
\begin{tabularx}{\textwidth}{@{}>{\raggedright\arraybackslash}p{2.35cm} XXXX@{}}
\toprule
\textbf{Aspecto} & \textbf{Excelente (40)} & \textbf{Bueno (30)} & \textbf{Aceptable (20)} & \textbf{Bajo (10)}\\
\midrule

\textbf{Reproducción y rigor experimental}
& Reproduce la l\'inea base, detecta y trata los duplicados, y todo resultado viene con semillas m\'ultiples y dispersi\'on. El cuaderno corre completo.
& Reproduce la l\'inea base con un protocolo correcto; la auditor\'ia del dato o la repetici\'on con semillas es parcial.
& Reproduce resultados pero con un protocolo discutible (fuga posible, una sola semilla, m\'etrica \'unica).
& No reproduce la l\'inea base o el cuaderno no corre; las cifras no son verificables.\\
\addlinespace[3pt]

\textbf{Diseño de los experimentos}
& Los barridos de T2--T5 est\'an bien dise\~nados, con rangos justificados y controles que permiten aislar la causa.
& Los experimentos son correctos pero alguno tiene rango insuficiente o le falta un control.
& Ejecuta los experimentos m\'inimos sin justificar las decisiones de dise\~no.
& Los experimentos no permiten concluir nada: rangos arbitrarios, sin controles, o incompletos.\\
\addlinespace[3pt]

\textbf{Interpretación teórica (cap.~20)}
& Usa los conceptos de la clase como instrumento de medici\'on; distingue con evidencia entre capacidad, optimizaci\'on y sesgo inductivo, y reconoce lo que sus datos \emph{no} permiten afirmar.
& Aplica los conceptos correctamente y defiende una tesis, con alguna afirmaci\'on que excede la evidencia.
& Menciona los conceptos correctos pero la conexi\'on con sus propios resultados es d\'ebil o decorativa.
& Usa el vocabulario de forma err\'onea o repite las conclusiones del art\'iculo sin contrastarlas.\\
\addlinespace[3pt]

\textbf{Análisis de datos y evidencia}
& An\'alisis profundo: figuras legibles y bien elegidas, incertidumbre reportada, comparaciones justas entre modelos.
& An\'alisis s\'olido; las figuras comunican, con mejoras posibles en claridad o en el tratamiento de la incertidumbre.
& An\'alisis adecuado pero superficial; faltan m\'etricas por clase o las figuras no aportan al argumento.
& An\'alisis superficial o inconsistente con las cifras del cuaderno.\\
\addlinespace[3pt]

\textbf{Comunicación del informe}
& Preciso, bien estructurado y sin errores; cada afirmaci\'on es rastreable hasta una figura o tabla.
& Bien organizado, con pocos errores menores.
& Se entiende, pero la estructura o la redacci\'on dificultan seguir el argumento.
& Desorganizado, impreciso o con errores graves que impiden la comprensi\'on.\\
\addlinespace[3pt]

\textbf{Lectura clínica y ética}
& Trata el costo asim\'etrico del error con una matriz justificada, reeval\'ua bajo ese criterio y discute con seriedad los l\'imites de transferir el modelo a otro sistema de salud.
& Aborda el costo asim\'etrico y los l\'imites del despliegue, con menos profundidad o sin reevaluar los modelos.
& Menciona el problema \'etico de forma gen\'erica, sin consecuencias sobre el an\'alisis.
& Omite la dimensi\'on cl\'inica o la reduce a una frase de cierre.\\
\bottomrule
\end{tabularx}

\normalsize
\vspace{1em}
\noindent\textbf{Conversión a nota} (60\,\% de exigencia, $P$ = puntaje sobre 240):
\[
\text{Nota} =
\begin{cases}
1{,}0 + 3{,}0\,\dfrac{P}{144} & \text{si } P < 144,\\[8pt]
4{,}0 + 3{,}0\,\dfrac{P-144}{96} & \text{si } P \geq 144.
\end{cases}
\]

\vspace{0.4em}
\noindent
\begin{tabularx}{\textwidth}{@{}l *{8}{>{\centering\arraybackslash}X}@{}}
\toprule
\textbf{Puntaje} & 48 & 96 & 120 & \textbf{144} & 168 & 192 & 216 & \textbf{240}\\
\textbf{Nota}    & 2{,}0 & 3{,}0 & 3{,}5 & \textbf{4{,}0} & 4{,}75 & 5{,}5 & 6{,}25 & \textbf{7{,}0}\\
\bottomrule
\end{tabularx}

% =====================================================================
\section{Referencias}

\begin{itemize}
	\item Tzimourta, K.\,D., Tsipouras, M.\,G., Angelidis, P., Tsalikakis, D.\,G. \& Orovou, E. (2025). \emph{Maternal Health Risk Detection: Advancing Midwifery with Artificial Intelligence}. Healthcare, 13(7), 833.\\ \url{https://doi.org/10.3390/healthcare13070833}
	\item Ahmed, M. (2023). \emph{Maternal Health Risk} [conjunto de datos]. UCI Machine Learning Repository.\\ \url{https://doi.org/10.24432/C5DP5D}
	\item Prince, S.\,J.\,D. \emph{Understanding Deep Learning}, cap\'itulo 20 (\emph{Why does deep learning work?}) y cap\'itulo 8 (doble descenso). MIT Press.\\ \url{https://udlbook.github.io/udlbook/}
	\item Zhang, C., Bengio, S., Hardt, M., Recht, B. \& Vinyals, O. (2017). \emph{Understanding deep learning requires rethinking generalization}. ICLR.
	\item Frankle, J. \& Carbin, M. (2019). \emph{The lottery ticket hypothesis: finding sparse, trainable neural networks}. ICLR.
	\item Keskar, N.\,S., Mudigere, D., Nocedal, J., Smelyanskiy, M. \& Tang, P.\,T.\,P. (2017). \emph{On large-batch training for deep learning: generalization gap and sharp minima}. ICLR.
	\item Grinsztajn, L., Oyallon, E. \& Varoquaux, G. (2022). \emph{Why do tree-based models still outperform deep learning on typical tabular data?} NeurIPS Datasets and Benchmarks.
\end{itemize}

\vspace{1em}
\hrule
\vspace{0.4em}
\noindent{\footnotesize\color{StanfordGray}
Material del curso MDS221, Mag\'ister en Data Science, Universidad Cat\'olica del Maule. El conjunto de datos se distribuye bajo licencia CC BY 4.0 y el art\'iculo de referencia bajo CC BY.}

\end{document}
